% 00-kurzfassung.tex — Kurzfassung (deutsch) und Abstract (englisch)
% Weniger Leerraum über der Überschrift, damit die Kurzfassung auf einer Seite bleibt
{\renewcommand*{\chapterheadstartvskip}{\vspace*{-1.5cm}}\addchap{Kurzfassung}}

Das Training bildbasierter Reinforcement-Learning-Agenten für die Fahrzeugregelung ist in realitätsnahen 3D-Simulatoren rechenintensiv: Ein Trainingslauf bindet Tage an Rechenzeit, bevor feststeht, ob Beobachtung, Belohnungsfunktion und Lernalgorithmus zueinander passen. Die vorliegende Arbeit begegnet diesem Problem am Beispiel der frei fahrenden Miniaturfahrzeuge des Systems \textit{Carrera Hybrid} (Maßstab etwa 1:43) in zwei Phasen. In Phase 1 werden die konzeptionellen Fragen in einer leichtgewichtigen 2D-Simulationsumgebung geklärt, in Phase 2 wird die dabei entwickelte Pipeline in die physikalisch simulierte, gerenderte 3D-Umgebung NVIDIA Isaac Sim übertragen.

In Phase 1 wurden fünf Beobachtungsvarianten mit den Algorithmen Proximal Policy Optimization (\ac{PPO}) und Soft Actor-Critic (\ac{SAC}) verglichen. Als methodischer Kern entstand eine Pipeline aus privilegierter Expertendestillation: Ein Lehrer-Agent mit exakter Beobachtung erzeugt Demonstrationsdaten, auf denen ein Faltungsnetz (\ac{CNN}) überwacht vortrainiert wird, bevor SAC den Kamera-Agenten online weiter optimiert. Das zentrale Ergebnis dieser Phase ist ein Ausschlusskriterium. Nach der Umstellung auf ein gemessenes Fahrzeugmodell konvergierten die Kamera-Agenten nicht mehr, während ein Kontrollversuch mit exakter Beobachtung dieselbe Aufgabe in \SI{54}{\minute} löste. Ursache ist eine Wahrnehmungsschwelle: Unterhalb von \SI{0.64}{\meter\per\second} verschiebt sich das Fahrzeug in der \num{166}\,$\times$\,\num{100}\,Pixel großen Beobachtung um weniger als einen Bildpunkt je Einzelbild, sodass der Bildstapel keine Geschwindigkeitsinformation trägt. Zudem erwies sich ein fester Startpunkt als Fehlerquelle, da er auswendig gelernte Trajektorien statt Fahrfähigkeit misst.

In Phase 2 wurde die Pipeline in Isaac Sim und Isaac Lab mit übernommener Längsdynamik, gerendertem Kamerabild und parallelen Umgebungen umgesetzt. Die vorab berechnete Wahrnehmungsschwelle liegt dort mit \SI{0.32}{\meter\per\second} unter dem Fahrtempo. Neu zeigte sich, dass das Reinforcement Learning einen vortrainierten Encoder zerstört, wenn dieser zu Beginn nicht eingefroren wird. Unter diesen Bedingungen entstanden zwei Kamera-Agenten, die zuverlässig Runden fahren. Der erste ist auf eine Darstellung der Strecke spezialisiert und nur rund \SI{4}{\percent} langsamer als ein Lehrer mit exaktem Zustand. Der zweite wurde auf zwölf Darstellungen mit Bildstörungen trainiert und fährt auf allen, auf den besten 19 Runden ohne Crash. Eine andere Streckenführung fährt keiner der beiden: Übertragbar ist das Sehen, nicht die Fahrweise.

Die Arbeit zeigt damit, dass ein Kamera-Agent unter realistischer Physik in der Simulation fahren lernen kann. Für den Transfer auf die reale Bahn werden die nächsten Schritte abgeleitet, vor allem das Training auf vielen Streckenführungen und die weitere Angleichung der Simulation an die Realität.

\clearpage

\addchap{Abstract}

Training vision-based reinforcement learning agents for vehicle control is computationally expensive in realistic 3D simulators: a single training run occupies days of compute before it becomes clear whether observation, reward function and learning algorithm fit together at all. This work addresses the problem in two phases, using the free-running miniature vehicles of the \textit{Carrera Hybrid} system (scale approximately 1:43). In phase 1, the conceptual questions are settled in a lightweight 2D simulation environment; in phase 2, the resulting pipeline is transferred to the physically simulated, rendered 3D environment NVIDIA Isaac Sim.

In phase 1, five observation variants were compared using Proximal Policy Optimization (\ac{PPO}) and Soft Actor-Critic (\ac{SAC}). The methodological core is a privileged expert distillation pipeline: a teacher agent with exact observations generates demonstration data on which a convolutional neural network (\ac{CNN}) is pretrained in a supervised fashion before SAC optimises the camera agent online. The central result of this phase is an exclusion criterion. After switching to a measured vehicle model, the camera agents stopped converging, while a control experiment with exact observations solved the same task in \SI{54}{\minute}. The cause is a perception threshold: below \SI{0.64}{\meter\per\second}, the vehicle moves by less than one pixel per frame in the \num{166}\,$\times$\,\num{100}\,pixel observation, so the frame stack carries no velocity information. In addition, a fixed starting position proved to be a source of error, as it measures memorised trajectories rather than driving ability.

In phase 2, the pipeline was implemented in Isaac Sim and Isaac Lab with the measured longitudinal dynamics, a rendered camera image and parallel environments. The perception threshold, computed in advance, lies at \SI{0.32}{\meter\per\second} and thus below the driving speed. A new finding is that reinforcement learning destroys a pretrained encoder unless it is frozen at the start of training. Under these conditions, two camera agents emerged that reliably complete laps. The first is specialised on a single appearance of the track and only about \SI{4}{\percent} slower than a teacher with exact state information. The second was trained on twelve appearances with image perturbations and drives on all of them, completing 19 laps without a crash on the best ones. Neither agent drives a different track layout: what transfers is perception, not driving behaviour.

The work thus shows that a camera agent can learn to drive under realistic physics in simulation. For the transfer to the real track, the next steps are derived, above all training on many track layouts and further aligning the simulation with reality.

\clearpage
